\documentclass[10pt,a4paper]{amsart}
\usepackage[margin=19mm]{geometry}
\usepackage{amsmath,amssymb,mathtools}
\usepackage[scr=rsfs,cal=euler]{mathalfa}
\usepackage{xfrac}
\usepackage[unicode,hidelinks,bookmarks=false]{hyperref}
\newcommand{\ie}{i.e.\ }
\newcommand{\cf}{cf.\ }
\newcommand{\resp}{respectively\ }
\newcommand{\Subsection}{\subsection}
\providecommand{\nobreakdash}{\nobreak\hbox{-}}
\setlength{\emergencystretch}{2em}
\multlinegap0pt
\usepackage{bm}
\usepackage{bbm}
\usepackage{dsfont}
\usepackage{xspace}
\usepackage{cancel}
\usepackage{mathtools}
\usepackage{amsthm}
\makeatletter
\@ifundefined{theorem}{\newtheorem{theorem}{Theorem}}{}
\@ifundefined{proposition}{\newtheorem{proposition}[theorem]{Proposition}}{}
\makeatother
\DeclareMathOperator{\Real}{Re}
\DeclareMathOperator{\Ric}{Ric}
\DeclareMathOperator{\tr}{tr}
\newcommand{\dbar}{\overline{\partial}}
\newcommand{\e}{\varepsilon}
\newcommand{\ip}[1]{\left\langle#1\right\rangle}
\newcommand{\oo}[1]{\overline{#1}}
\newcommand{\parabolic}[1]{\left(\frac{\d}{\d t} - \Delta_{#1}\right)}
\renewcommand{\d}{\partial}
\title[Fano Fibrations and Twisted K\"ahler-Einstein Metrics II: Th]{Fano Fibrations and Twisted K\"ahler-Einstein Metrics II: The K\"ahler-Ricci Flow}
\author{Alexander Bednarek}
\date{Source: arXiv:2512.21910v1 (CC BY 4.0)}
\begin{document}
\maketitle

\noindent\emph{Source and licence.} Fano Fibrations and Twisted K\"ahler-Einstein Metrics II: The K\"ahler-Ricci Flow. Version arXiv:2512.21910v1 (CC BY 4.0), distributed under the Creative Commons Attribution 4.0 International licence, as recorded in the arXiv licence metadata for this exact version and shown on the abstract page https://arxiv.org/abs/2512.21910v1. Full rights notice: licenses/arxiv-2512.21910-CC-BY-4.0.txt.\par\smallskip

\noindent\textit{Editorial scope.} Counted unit: the Proposition whose source label is \texttt{None} in the article (source0.tex lines 1295--1304, source label \texttt{None}), delivered together with the article's own complete original proof of it (source0.tex lines 1306--1391, one \texttt{proof} environment of 86 lines). No mathematics is written, repaired or re-derived: statement and proof are the article's own text line for line. Required context retained uncounted: KRF (lines 215--217); VFC (lines 307--310); SKE (lines 719--721); norm\_grad\_v\_proof\_inequality (lines 1202--1208). Every definition, display and label of those lines is retained unchanged. Omitted from this package and not redistributed: 1--214, 218--306, 311--718, 722--1201, 1209--1294, 1305--1305, 1392--1605. Nothing inside a retained excerpt is omitted or altered: every retained body is delivered exactly as the article's lines are written, so no omission receipt of any kind accompanies this package. Citations: the retained excerpts carry no citation command at all, so no bibliography entry of the article is transcribed and no reference is redistributed. Reference adaptation: none was needed and none was made; every reference inside the delivered unit resolves inside it, so no unresolved reference or citation is produced. Printed slips kept literal: no printed word, symbol, hypothesis or number of the counted statement or of its proof was altered. Layout and class: the article's own class is \texttt{article} with packages including inputenc, geometry, amsmath, amssymb, amsthm, hyperref, tikz-cd; the wrapper is the lane's pinned \texttt{amsart} editorial wrapper at 10pt on A4, which restates the theorem-like environments the retained text needs and adds the article's own macro definitions that the retained text needs (\texttt{\textbackslash Real}, \texttt{\textbackslash Ric}, \texttt{\textbackslash d}, \texttt{\textbackslash dbar}, \texttt{\textbackslash e}, \texttt{\textbackslash ip}, \texttt{\textbackslash oo}, \texttt{\textbackslash parabolic}, \texttt{\textbackslash tr}), with extra wrapper packages (bm, bbm, dsfont, xspace, cancel, mathtools). The delivered numbering is the unit's own. Assets: no retained span contains a figure environment, a graphics-inclusion command, a PDF-page inclusion, a table or a TikZ picture; no image file is copied into this package, the package declares no asset dependency and no image file is redistributed with it. Licence: version arXiv:2512.21910v1 of this article is distributed under the Creative Commons Attribution 4.0 International licence, as recorded in the arXiv licence metadata for this exact version and shown on the abstract page https://arxiv.org/abs/2512.21910v1. A full rights notice is delivered at licenses/arxiv-2512.21910-CC-BY-4.0.txt. Characters: every retained excerpt is pure ASCII, so the delivered engine, XeLaTeX with its default Latin Modern text font, reports no missing character.
\begin{equation}\label{KRF}\tag{KRF}
    \frac{\d \omega(t)}{\d t} = -\Ric(\omega(t)) - \omega(t), \quad \omega(0) = \omega_0.
\end{equation}

\begin{align}\tag{VFC}\label{VFC}
    C_K^{-1}E(t)^{n-m}\Omega \leq \omega(t)^{n} \leq C_KE(t)^{n-m}\Omega \quad \text{on} \quad K \times [0,T),\\
    \omega(t)^n \leq U_XE(t)^{n-m} \Omega \quad \text{on} \quad X \times [0,T).\notag
\end{align}

\begin{equation}\tag{SKE}\label{SKE}
    \Ric \omega_{SKE,b} = \lambda \omega_{SKE,b}.
\end{equation}

\begin{align}\label{norm_grad_v_proof_inequality}
    \parabolic{}&\left(\frac{|\nabla v|^2}{v} + (1-e^{-T})\tr_{\omega(t)} f^*\omega'_B\right)\\\notag
    &\leq \frac{|\nabla v|^2 - |\nabla \nabla v|^2 - |\nabla \oo{\nabla} v|^2}{v}- \frac{|\nabla v|^2}{v^2} \left(m-(1-e^{-T})\tr_{\omega(t)} f^*\omega'_B - \frac{2A e^{-t}}{1-e^{-T}}\right)\\\notag
    &+ 2 \Real \ip{\nabla \left(\frac{|\nabla v|^2}{v}+ (1-e^{-T})\tr_{\omega(t)} f^*\omega'_B \right), \frac{\oo{\nabla} v}{v}}\\\notag
    &- 4 \Real \ip{\nabla (1-e^{-T})\tr_{\omega(t)}f^*\omega_B', \frac{\oo{\nabla} v}{v}}\\\notag
    &+ C - C^{-1}|\nabla (1-e^{-T})\tr_{\omega(t)} f^* \omega'_B|^2.\notag
\end{align}

\begin{proposition}
Assume that the (\ref{KRF}) satisfies (\ref{VFC}). If $\rho_{SPR} = 0$ on $X \backslash S$ or (\ref{SKE}) holds, then for any compact set $K \subseteq X \backslash S$ where $h_K$ is Lipschitz there exist a constant $C = C(n,m,T,\omega_0,f^*\eta,\rho_{SPR},f^*\rho_B',K,C_K,U_X) > 0$ or a constant $C = C(n,m,T,\omega_0,f^*\eta,\rho_{SKE},f^*\rho_B',K,C_K,U_X) > 0$, respectively, such that on $K \times [0,T)$,
\begin{equation*}
    \Delta_{\omega(t)} v \leq C.
\end{equation*}
On the other hand, assume that $f:X \to B$ is a submersion and (\ref{KRF}) satisfies (\ref{VFC}). If $\omega_{SPR} = \omega_{0}$ on $X$ or if (\ref{SKE}) holds, then there exists a $C = C(n,m,T,\omega_0, f^*\eta,\rho_{SPR},f^*\rho_B',C_X) > 0$ or $C = C(n,m,T,\omega_0,f^*\eta,\rho_{SKE},f^*\rho_B',C_X) > 0$, respectively, such that on $X \times [0,T)$ we have
\begin{equation*}
    \Delta_{\omega(t)} v \leq C.
\end{equation*}
\end{proposition}

\begin{proof}
Again, we only provide the argument for the first case where $\rho_{SPR} = 0$ on $X \backslash S$ and take $K \subset \subset U \subset \subset K' \subset \subset X \backslash S$. We can calculate the evolution equation of $\Delta v - (1-e^{-T})\tr_{\omega(t)} f^*\omega'_B$ over $X \backslash S \times [0,T)$ to be 
\begin{align*}
    \parabolic{}&(\Delta v - (1-e^{-T})\tr_{\omega(t)} f^*\omega'_B)\\ &= \Delta v - (1-e^{-T})\tr_{\omega(t)}f^*\omega'_B+ \ip{\Ric, -(1-e^{-T})f^* \omega'_B + i\d\dbar v}\\
    &= \Delta v - (1-e^{-T})\tr_{\omega(t)} f^*\omega'_B+\frac{e^{t-T}}{1-e^{t-T}} (\Delta v - (1-e^{-T})\tr_{\omega(t)} f^*\omega'_B )\\
    &+ \frac{1}{1-e^{t-T}}\left(|(1-e^{-T})f^*\omega'_B|^2 + |\nabla \oo{\nabla} v|^2 + 2 \Real \ip{(1-e^{-T})f^*\omega'_B, i\d\dbar v }\right).
\end{align*}
where the last line uses
\begin{equation*}
    (1-e^{t-T}) \Ric = e^{t-T} \omega(t) - (1-e^{-T})f^*\omega'_B + i\d\dbar v.
\end{equation*}
It follows that on $U \times [0,T)$ we have 
\begin{align*}
    \parabolic{} &\left(\frac{(1-e^{t-T})(\Delta v - (1-e^{-T})\tr_{\omega(t)} f^*\omega'_B)}{v}\right)\\
    &= (1-e^{t-T}) \left(\frac{\parabolic{} (\Delta v - (1-e^{-T})\tr_{\omega(t)} f^*\omega'_B)}{v}\right)\\
    &-e^{t-T} \frac{\Delta v - (1-e^{-T}) \tr_{\omega(t)}f^*\omega'_B}{ v}\\
    &- (1-e^{t-T})\left(\frac{\Delta v - (1-e^{-T})\tr_{\omega(t)} f^*\omega'_B}{v} \frac{\parabolic{} v}{v}\right)\\
    &+ 2\Real \ip{\nabla \left(\frac{(1-e^{t-T})(\Delta v - (1-e^{-T})\tr_{\omega(t)} f^*\omega'_B)}{v}\right), \frac{\oo{\nabla} v}{v}}\\
    &= (1-e^{t-T}) \frac{\Delta v - (1-e^{-T})\tr_{\omega(t)} f^*\omega'_B}{v} \left(1+\frac{-m+(1-e^{-T})\tr_{\omega(t)}f^*\omega'_B + \frac{2A e^{-t}}{1-e^{-T}}}{v}\right)\\
    &+\frac{|(1-e^{-T})f^*\omega'_B|^2 + |\nabla \oo{\nabla} v|^2 + 2 \Real \ip{(1-e^{-T})f^*\omega'_B, i\d\dbar v }}{v}\\
    &+ 2\Real \ip{\nabla \left(\frac{(1-e^{t-T})(\Delta v - (1-e^{-T})\tr_{\omega(t)} f^*\omega'_B)}{v}\right), \frac{\oo{\nabla} v}{v}}\\
    &\leq \frac{C}{v} + (1-e^{t-T}) \frac{\Delta v }{v} \left(1+\frac{-m+(1-e^{-T})\tr_{\omega(t)}f^*\omega'_B + \frac{2A e^{-t}}{1-e^{-T}}}{v}\right)+\frac{2|\nabla \oo{\nabla} v|^2}{v}\\
    &+ 2\Real \ip{\nabla \left(\frac{(1-e^{t-T})(\Delta v - (1-e^{-T})\tr_{\omega(t)} f^*\omega'_B)}{v}\right), \frac{\oo{\nabla} v}{v}}
\end{align*}
where we apply the bounds $|f^*\omega_B'|_{\omega(t)}^2 \leq (\tr_{\omega(t)}f^*\omega_B')^2$ and $0 \leq \tr_{\omega(t)} f^*\omega'_B \leq C$, on $K' \times [0,T)$. Recall inequality (\ref{norm_grad_v_proof_inequality}) from the previous Proposition so that over $U \times [0,T)$ we have
\begin{align*}
    \parabolic{}&\left(\frac{|\nabla v|^2}{v} + (1-e^{-T})\tr_{\omega(t)} f^*\omega'_B\right)\\
    &\leq \frac{|\nabla v|^2 - |\nabla \nabla v|^2 - |\nabla \oo{\nabla} v|^2}{v}+ C\frac{|\nabla v|^2}{v^2} + C\\
    &+ 2 \Real \ip{\nabla \left(\frac{|\nabla v|^2}{v}+ (1-e^{-T})\tr_{\omega(t)} f^*\omega'_B \right), \frac{\oo{\nabla} v}{v}},
\end{align*}
so that after setting 
\begin{equation*}
    Q_4 = \frac{(1-e^{t-T})(\Delta v - (1-e^{-T})\tr_{\omega(t)} f^*\omega'_B)}{v}, \quad Q_5 = Q_4 + 3 Q_3,
\end{equation*}
we have
\begin{align*}
    \parabolic{}Q_5 &\leq C + \frac{C}{v}-\frac{|\nabla \oo{\nabla} v|^2}{v} + 2\Real \ip{\nabla Q_5, \frac{\oo{\nabla} v}{v}}\\
    &+(1-e^{t-T}) \frac{\Delta v }{v} \left(1 + \frac{-m+(1-e^{-T})\tr_{\omega(t)}f^*\omega'_B + \frac{2A e^{-t}}{1-e^{-T}}}{v}\right)
\end{align*}
on $U \times [0,T)$. Next, observe that $n|\nabla \oo{\nabla} v|^2 \geq (\Delta v)^2$ so that
\begin{align*}
    \parabolic{} Q_5 &\leq \frac{C}{v} -\frac{\frac{1}{n}(\Delta v)^2}{v} + 2\Real \ip{\nabla Q_5, \frac{\oo{\nabla} v}{v}}\\
    &+ (1-e^{t-T}) \frac{\Delta v }{v} \left(1+\frac{-m+(1-e^{-T})\tr_{\omega(t)}f^*\omega'_B + \frac{2A e^{-t}}{1-e^{-T}}}{v}\right)\\
    &\leq \frac{C}{v} + \frac{(\e-\frac{1}{n})(\Delta v)^2}{v} + 2 \Real \ip{\nabla Q_5, \frac{\oo{\nabla} v}{v}}
\end{align*}
for some $\e > 0$ to be fixed. As before, taking $\Upsilon = e^{-C_1\chi^{-C_2}}$, we find that on $X \times [0,T)$,
\begin{equation*}
    \parabolic{} \Upsilon Q_5 = \Upsilon \parabolic{} Q_5 + \left(\frac{|\nabla \Upsilon|^2}{\Upsilon} - \Delta \Upsilon\right) Q_5 - 2 \Real \ip{\nabla(\Upsilon Q_5), \frac{\oo{\nabla} \Upsilon}{\Upsilon}}.
\end{equation*}
We estimate the middle term by
\begin{equation*}
    \left(\frac{|\nabla \Upsilon|^2}{\Upsilon} - \Delta \Upsilon\right) Q_5 \leq C\chi^{-2C_2-2}\Upsilon |Q_5| \leq C\chi^{-4C_2-4}\Upsilon + \frac{\e(\Delta v)^2}{v}\Upsilon.
\end{equation*}
We also note that within the first term we need the following estimate
\begin{align*}
    2 \Upsilon \Real \ip{\nabla Q_5, \frac{\oo{\nabla}v}{v}} &= 2 \Real \ip{\nabla (\Upsilon Q_5), \frac{\oo{\nabla}v}{v}} - 2Q_5 \Real \ip{\nabla \Upsilon, \frac{\oo{\nabla}v}{v}}\\
    &\leq 2 \Real \ip{\nabla (\Upsilon Q_5), \frac{\oo{\nabla}v}{v}} + \frac{C|Q_5|}{\sqrt{v}} \chi^{-2C_2-2} \Upsilon\\
    &\leq 2 \Real \ip{\nabla(\Upsilon Q_5), \frac{\oo{\nabla}v}{v}} + C\chi^{-4C_2-4}\Upsilon + \frac{\e(\Delta v)^2}{v}\Upsilon.
\end{align*}
All together, we currently have
\begin{align*}
    \parabolic{} \Upsilon Q_5 &\leq C\chi^{-4C_2-4}\Upsilon+ \frac{C}{v}\Upsilon + \frac{(3\e-\frac{1}{n})(\Delta v)^2}{v}\Upsilon\\
    &+ 2 \Real \ip{\nabla(\Upsilon Q_5), \frac{\oo{\nabla}v}{v}}-2 \Real \ip{\nabla(\Upsilon Q_5), \frac{\nabla \Upsilon}{\Upsilon}}.
\end{align*}
Now fix $\e = \frac{1}{6n}$ so that $3\e - \frac{1}{n} = -\frac{1}{2n}$. Assume that $(x^*,t^*) \in X \times [0,T)$ is a point at which $\Upsilon Q_5$ achieves a maximum, if $x^* \in X \backslash U$ then $(\Upsilon Q_5)(x^*,t^*) = 0$ is bounded. Otherwise, $x^* \in U$ and
\begin{equation*}
    0 \leq  C \chi(x^*)^{-4C_2-4}\Upsilon(x^*)+ \frac{C}{v(x^*,t^*)}\Upsilon(x^*) - \frac{\frac{1}{2n}(\Delta v)(x^*,t^*)^2}{v(x^*,t^*)}\Upsilon(x^*)
\end{equation*}
so that
\begin{equation*}
    (\Delta v)(x^*,t^*) \leq C
\end{equation*}
since $0 < v(x^*,t^*) \leq C$, and we see
\begin{equation*}
    (\Upsilon Q_5)(x^*,t^*) \leq C.
\end{equation*}
It follows that for all $(x,t) \in X \times [0,T)$ we have
\begin{equation*}
    (\Upsilon Q_5)(x,t) \leq C
\end{equation*}
so that $\Upsilon Q_4 \leq C$ and thus
\begin{equation*}
    \Delta v \leq \frac{C\Upsilon^{-1}v}{(1-e^{t-T})} + (1-e^{-T})\tr_{\omega(t)} f^*\omega'_B.
\end{equation*}
The result follows.
\end{proof}

\end{document}
